% acmart conversion deferred until venue choice — swap documentclass + \cite style then.
\documentclass[11pt]{article}

\usepackage[utf8]{inputenc}
\usepackage[T1]{fontenc}
\usepackage{lmodern}
\usepackage{amsmath,amssymb,amsthm}
\usepackage{mathtools}
\usepackage{mathpartir}
\usepackage{booktabs}
\usepackage{graphicx}
\usepackage{geometry}
\usepackage{xcolor}
\usepackage[activate=false]{microtype}
\usepackage[hidelinks]{hyperref}
\usepackage{listings}

\geometry{margin=1.25in}

\lstset{
  basicstyle=\small\ttfamily,
  breaklines=true,
  columns=flexible,
  mathescape=true,
  keepspaces=true,
}

%% ============================================================
%%  Notation (shared with the till paper where the fragments overlap)
%% ============================================================
\newcommand{\lax}[2]{\{#1\}\mathbin{@}#2}       % graded lax monad {S}@d
\newcommand{\laxb}[1]{\{#1\}}                   % bare lax monad {S}
\newcommand{\gbang}[1]{{!}_{#1}}                 % counted bang  !_k
\newcommand{\loli}{\multimap}                    % linear implication ⊸
\newcommand{\tensor}{\otimes}                    % tensor ⊗
\newcommand{\addconj}{\mathbin{\&}}              % additive conjunction &
\newcommand{\adddisj}{\oplus}                    % additive disjunction ⊕
\newcommand{\bang}{{!}}                          % bang !
\newcommand{\QQ}{\mathbb{Q}}                     % rationals
\newcommand{\NN}{\mathbb{N}}                     % naturals
\newcommand{\CC}{\mathbb{C}}                     % complex numbers
\newcommand{\stmp}[2]{#1\mathbin{@}#2}           % stamp  a@t
\newcommand{\enc}[1]{\lceil #1 \rceil}          % state encoding bracket
\newcommand{\till}{\textsc{till}}                % the time instance
\newcommand{\gill}{\textsc{gill}}                % the graded family
\newcommand{\will}{\textsc{will}}                % the measure instance (this paper)
\newcommand{\calc}{\textsc{calc}}                % sandbox name

%% will-specific
\newcommand{\erho}[2]{\exists_{\rho}\, #1{:}#2.\;} % weighted existential ∃_ρ x:s.
\newcommand{\dtok}[2]{\mathsf{drawn}\;#1\;#2}      % draw token  drawn c s
\newcommand{\trz}[1]{\langle #1\rangle}            % token zone ⟨Θ⟩
\newcommand{\Tr}{\Theta}                           % a trace
\newcommand{\ghost}{\mathsf{ghost}}
\newcommand{\focR}{\mathrel{\gg}}                  % right focus  Δ ≫ P
\newcommand{\invR}{\mathrel{\Uparrow}}             % inversion
\newcommand{\seq}[3]{#1;\;#2\vdash #3}             % Γ; Δ ⊢ A
\newcommand{\seqt}[4]{#1;\;#2\vdash #3\;[#4]}      % Γ; Δ ⊢ A [Θ]
\newcommand{\wt}[1]{w(#1)}

%% Theorem environments — numbered per section
\newtheorem{theorem}{Theorem}[section]
\newtheorem{lemma}[theorem]{Lemma}
\newtheorem{corollary}[theorem]{Corollary}
\newtheorem{proposition}[theorem]{Proposition}
\theoremstyle{definition}
\newtheorem{definition}[theorem]{Definition}
\theoremstyle{remark}
\newtheorem{remark}[theorem]{Remark}

\title{Weight in the Endsequent:\\
  Cut Admissibility for a Measure-Weighted Existential\\
  by Draw-Token Internalization}
\author{Denis Erfurt}
\date{Draft (TODO\_0283)}

\begin{document}
\maketitle

\begin{abstract}
We present \will{}, an intuitionistic linear sequent calculus with a
\emph{measure-weighted existential quantifier} $\erho{x}{s} A$: the
binder ranges over a closed classifier sort $s$ whose constructors carry
declared prior weights $\rho : s \to \QQ_{\geq 0}$, and a derivation's
weight is the product of the priors of the witnesses it draws.  The
central difficulty is that weight, treated as a judgment grade, is
\emph{not preserved by cut elimination}: the principal reduction of the
weighted existential substitutes the witness and deletes the
introduction step, losing its factor.  Our resolution is to stop
annotating: the run trace is \emph{internalized} as a multiset of
linear \emph{draw-token} hypotheses $\dtok{c}{s}$ in the ordinary
linear zone.  The weighted right rule \emph{consumes} a token; weight
becomes a function of the \emph{endsequent}; and since cut elimination
preserves endsequents by construction, exact weight conservation is
definitional.  We prove cut admissibility for the extended calculus
(with no relaxation order, no expectation, no threshold --- raw
$\QQ_{\geq0}$ priors conserved on the nose), and derive three
structural consequences: \emph{ghost} weakening restricted to the token
class is forced by closure under cut reduction; \emph{no promotion
through a draw} (the $\bang$R empty-zone condition makes duplication
and erasure trace-safe); and \emph{identity expansion fails} at
$\exists_\rho$ --- a proof-theoretic no-cloning property, which forces
the focused system to treat $\exists_\rho$ as a synthetic atom.  We
give the focused system and the permutation argument for its
completeness, define the \emph{proof-counting provenance polynomial}
$P_A \in \NN[X]$ whose evaluation homomorphisms recover the
probabilistic mass ($\QQ_{\geq0}$), the tropical optimum
($\min,+$), and a reserved amplitude face ($\CC$), and prove that the
operational \emph{decimation driver} is adequate: its collapse paths
are in bijection with the focused ghost-free derivations the
polynomial counts.  The calculus is implemented: every theorem in the
paper is either executable or checked by an LF-style kernel, and
complete runs of the stochastic engine ship as kernel-verified
certificates whose endsequents carry their traces.
\end{abstract}

\tableofcontents
\bigskip\hrule\bigskip

%%====================================================================
\section{Introduction}
\label{sec:intro}
%%====================================================================

\subsection{The problem}

Linear logic is the natural home for stochastic \emph{generation}: a
wave-function-collapse tiler, a PCFG sampler, and a probabilistic
multiset rewriting engine all manipulate resources that are produced,
consumed, and --- crucially --- \emph{drawn}.  The quantitative
refinements of linear logic, however, live almost entirely on the
\emph{necessity} side: graded exponentials
$\gbang{r}A$~\cite{GirardScedrovScott1992,Atkey2018,OrchardLiepeltEades2019},
graded monads as categorical
structure~\cite{Katsumata2014,FujiiKatsumataMillies2016}, and (in the
predecessor to this paper) a delay-graded lax modality with cut
admissibility~\cite{TillPaper}.  What generation needs is different: a
\emph{weighted existential}.  ``There exists a tile for this cell,
distributed by the priors $\rho$'' is a quantifier, not a modality; its
witness is chosen by a draw; and the weight of a run is the product of
the priors of its draws.

The na\"ive design annotates the judgment:
$\seq{\Gamma}{\Delta}{A\,\langle w\rangle}$, with the weighted right
rule multiplying $\rho(c)$ into $w$.  This design is broken, and the
break is instructive.  The principal cut of the weighted existential
against its left rule reduces by \emph{substituting the witness and
deleting the introduction step} --- and the deleted step is the only
place the factor $\rho(c)$ entered.  On a two-constructor sort the
annotated calculus rewrites a weight-$2$ derivation to a weight-$1$
derivation of the same sequent: cut elimination is not
weight-preserving, so the annotation does not denote anything stable
under the equational theory of proofs.  The same wall --- a sum- or
measure-valued judgment without syntactic cut elimination --- appears
independently in counting propositional logic~\cite{AntonelliDLP2022},
whose semantics is exact but whose cut rule is only semantically
admissible, and in the algebraic $\lambda$-calculus~\cite{Vaux2009},
where scalar sums survive reduction only under positivity constraints.
The only syntactic success we know of, cut elimination for Riesz modal
logic~\cite{LucasMio2022}, pays with hypersequents, a scalar-weighted
context calculus, and a fifty-page proof --- and its interaction with
linear multiplicatives is open.

\subsection{The resolution: internalize the trace}

Our resolution is to stop annotating.  A stochastic run already has a
canonical record: its \emph{trace}, the multiset of draws it made.  We
internalize the trace into the sequent itself as ordinary linear
hypotheses.  One kernel-reserved atom family
\[
  \dtok{c}{s} \qquad (c \text{ a constructor of the classifier sort } s)
\]
is added, and the trace judgment is \emph{defined}, not axiomatized:
\[
  \seqt{\Gamma}{\Delta}{A}{\Tr} \;\;:=\;\;
  \seq{\Gamma}{\Delta,\trz{\Tr}}{A},
\]
where $\trz{\Tr}$ is the multiset of tokens $\dtok{c}{s}$ for the
records in $\Tr$.  The weighted right rule \emph{consumes} a token ---
``given that a draw of $c$ occurred, assert the existential via $c$''
--- and the weight of a sequent is a \emph{function of its endsequent}:
$\wt{\Tr} = \prod_{(c:s)\in\Tr}\rho(c)$.  Cut elimination preserves
endsequents by construction; conservation of the trace, and hence of
the weight, is \emph{definitional}.  The proof burden moves to where it
belongs: showing that the required reduction steps \emph{exist} in the
extended calculus.  They do (\S\ref{sec:cut}), and the three places
where existence is subtle each yield a structural theorem rather than a
patch:

\begin{itemize}
\item The principal reduction turns a \emph{used} draw into an
  \emph{unused} one, so the calculus needs \emph{ghost}: weakening
  restricted to the token class.  Tape-linear systems --- where tokens
  must be consumed exactly --- are \emph{not closed under cut
  reduction} (\S\ref{sec:cut}, case (a)).  Ghost is then excluded from
  the canonical counting class exactly the way ordinary weakening is
  excluded under focusing (\S\ref{sec:ghost}).
\item Duplication and erasure of subderivations, the classical dangers
  for any resource-counting discipline, are governed by one existing
  condition: promotion ($\bang$R) requires an \emph{empty linear zone}
  --- which now includes tokens.  \textbf{No promotion through a draw}:
  a derivation that draws is not promotable, no draw is ever duplicated
  or erased (\S\ref{sec:cut}, cases (e), (f)).  Probabilistic
  conclusions are not bangable knowledge; one may bang the
  \emph{consequences of an outcome}, never the luck itself.
\item Identity expansion \emph{fails} at $\exists_\rho$: eta-expanding
  $\erho{x}{s}B \vdash \erho{x}{s}B$ would need a token the sequent
  does not have.  A committed draw can be passed along whole but not
  deconstructed and re-derived --- re-derivation would be a second
  draw, a different endsequent (\S\ref{sec:nocloning}).  A
  no-cloning-flavored theorem appears at exactly the connective that
  carries randomness, in a classical calculus, for free; its design
  consequence is that the focused system must treat $\exists_\rho$ as a
  \emph{synthetic atom} on the identity dimension (\S\ref{sec:focused}).
\end{itemize}

\subsection{From counting to measure: provenance}

With cut admissibility in hand, the quantitative story is a
\emph{counting} story.  For a token-free boundary judgment we define
$N(A,\Tr)$, the number of focused, cut-free, ghost-free derivations of
$\seq{\Gamma}{\Delta,\trz{\Tr}}{A}$, and package the counts as a
\emph{proof-counting provenance polynomial}
$P_A = \sum_\Tr N(A,\Tr)\cdot X^{\Tr} \in \NN[X]$.  Because tokens are
linear, each derivation carries exactly \emph{one} monomial --- its
endsequent's token multiset --- sharpening the provenance-semiring
picture~\cite{GreenKT2007,GradelTannen2024} in which monomials annotate
proof trees.  By $\NN[X]$-initiality every semiring-valued reading is
an evaluation homomorphism out of $P_A$: at
$K=(\QQ_{\geq0},+,\cdot)$, $X_c \mapsto \rho(c)$, one obtains \will{}'s
probabilistic mass; at $K = (\min,+)$ the sum collapses onto the
single best derivation --- which is exactly why the
cost-graded sibling calculus (\gill{}) needs no trace and a per-proof
grade suffices; $K=\CC$ is reserved: evaluation is still a
homomorphism, but interference is forgetting-by-cancellation that no
trace-level statement can undo.  Cut admissibility is proved once; each
quantitative face inherits invariance through its homomorphism
(\S\ref{sec:bridge}).

\subsection{The engine, certified}

\will{} is not a paper calculus.  Its operational face is a
\emph{decimation driver} over committed-choice multiset rewriting: the
weighted existential in a rule consequent suspends as a \emph{wave};
the driver opens waves with fresh existential variables, re-settles so
that derived \emph{bias} facts condition the posterior, draws
minimum-entropy-first, and substitutes witnesses.  We prove the driver
\emph{adequate}: its collapse paths are in bijection with the focused
ghost-free derivations counted by $N$ (\S\ref{sec:adequacy}).  The
adequacy is not just stated but \emph{shipped}: a complete stochastic
run elaborates --- post-hoc grounded --- into a single derivation tree
checked by an LF-style kernel, whose endsequent
$\Gamma;\Delta_0,\trz{\Tr}\vdash\laxb{\bigotimes\text{residual}}$
carries the trace, so the run's prior mass is readable off the
certificate's endsequent alone, and a doctored weight or a truncated
token zone fails verification (\S\ref{sec:impl}).

\subsection{Contributions}

\begin{enumerate}
\item \textbf{The trace judgment by internalization}
  (\S\ref{sec:calculus}): the first cut-admissibility result for a
  measure-weighted existential in a linear sequent calculus, with
  \emph{exact} weight conservation.  The mechanism --- weight as a
  function of which hypotheses remain in the endsequent --- is new at
  the theorem level; \S\ref{sec:related} distinguishes the idea-level
  relatives.
\item \textbf{Three structural theorems} (\S\ref{sec:cut}--\ref{sec:ghost}):
  ghost forced by closure under reduction; no promotion through a draw;
  no-cloning as identity-expansion failure.
\item \textbf{The focused counting system} (\S\ref{sec:focused}): a
  focused calculus in which ghost-freeness is structural, $\exists_\rho$
  is a synthetic atom, and the completeness argument's new permutation
  cases are worked in full.
\item \textbf{Universality} (\S\ref{sec:bridge}): the proof-counting
  provenance polynomial and its evaluation faces; one cut theorem, all
  faces inherit.
\item \textbf{Driver adequacy and a certified engine}
  (\S\ref{sec:adequacy}--\ref{sec:impl}): collapse paths $\cong$
  counted derivations; complete stochastic runs as kernel-checked
  certificates whose endsequents carry their traces; a
  wave-function-collapse case study whose exact-mode measure matches
  brute-force enumeration.
\end{enumerate}

\subsection{Non-claims}

All proofs are on paper, none machine-checked; the kernel that checks
\emph{derivations} is an implementation artifact, not a proof
assistant, and mechanizing \S\ref{sec:cut} is future work.  The
focalization theorem's base cases (MALL focusing) are standard and
cited, not re-proved; \S\ref{sec:focused} works the cases our
connectives add.  Driver adequacy is scoped: boundary judgments,
conflict-free settle segments, and a fixed draw policy (the
policy-independence of the resulting measure follows, but scheduling
fairness for non-confluent programs is out of scope).  Convergence for
recursive sorts (subcritical priors), the unbiasedness of the
importance-weighted sampler, and the exactness of the underlying timed
engine are imported from the companion theory
documents~\cite{THY0026,TillPaper} and used as black boxes.  The
$\CC$-valued face is reserved, not developed.  A mass \emph{judgment}
(sequents proving ``the total mass of $A$ is $m$'') is deliberately
absent: the literature's price tags for syntactic cut elimination at
that level~\cite{LucasMio2022,AntonelliDLP2022} are the reason this
paper counts proofs instead.

%%====================================================================
\section{The calculus}
\label{sec:calculus}
%%====================================================================

\subsection{Base: a graded linear sequent calculus}

\will{} extends \gill{}, the graded intuitionistic linear sequent
calculus of the companion paper~\cite{TillPaper}: two-zone judgments
$\seq{\Gamma}{\Delta}{A}$ (persistent; linear), multiplicatives
$\tensor, 1, \loli$, additive $\addconj$, the exponential $\bang A$
with an explicit promotion rule requiring an empty linear zone, counted
parcels $\gbang{k}A$, a delay-graded lax monad $\lax{S}{d}$ whose side
conditions are theory premises over decidable arithmetic, and stamped
atoms $\stmp{a}{t}$.  Everything in this paper is orthogonal to the
grading of the monad: delay grades live on formulas, weight will live
in the zone, and the two never share a slot (\S\ref{sec:cut}, case
(g)).  We write the fragment's rules only as needed; the full rule set
is in the artifact.

Two base-calculus facts are load-bearing and worth stating now:

\begin{itemize}
\item \textbf{Promotion has an empty linear zone.}
  $\bang$R (and its counted variant at count $0$) concludes from
  $\seq{\Gamma}{\cdot}{A}$.  During this work an adversarial audit of
  the counted rule found it as originally written with an
  \emph{arbitrary} linear context --- an affine leak deriving
  $a \vdash \gbang{0}b$, harmless in the base calculus but fatal for
  trace conservation once affine tokens exist.  The repair (empty zone,
  like $1$R) restores $\gbang 0 A \dashv\vdash 1$ and is assumed
  throughout (\S\ref{sec:cut}, case (f$'$)).
\item \textbf{Identity is primitive at every formula.}  The base
  calculus takes $A \vdash A$ as a rule for all $A$; identity
  \emph{expansion} (deriving compound id from atomic id) is a theorem
  there.  In \will{} that theorem acquires an exception
  (\S\ref{sec:nocloning}), which is why primitivity matters.
\end{itemize}

\subsection{Sorts, priors, and tokens}

A \emph{classifier sort} $s$ is a closed-world sort whose members are
declared constructors; each member $c$ may carry a declared prior
weight $\rho(c) \in \QQ_{\geq0}$ (unannotated members default to $1$;
weights are unnormalized --- normalization is a boundary operation, as
in~\cite{Staton2017}).  Members may be nullary (\emph{rung~1}) or
constructors with classifier-sorted arguments (\emph{rung~2}), in which
case draws instantiate one head at a time and argument positions become
new draws --- ancestral PCFG generation; convergence is governed by
Chi--Geman subcriticality~\cite{ChiGeman1998} and imported
from~\cite{THY0026}.

The one new atom family is the \textbf{draw token}:
\[
  \dtok{c}{s}\colon \text{linear, ground, kernel-reserved.}
\]
\emph{Fences:} (i) no program rule may produce or match a token, no
right rule exists for it, and the forward-engine checker never
concludes one --- tokens enter judgments only at the boundary, minted
by the draw checker of \S\ref{sec:impl}, one per collapse event,
with $\rho(c)$ checked against the declared prior; (ii) tokens are
ground --- a draw records a constructor \emph{head}, so substitution
never touches tokens (rung-2 lazy draws record the head only, matching
priors being per-constructor).

\begin{definition}[Trace judgment]\label{def:trace}
For a multiset $\Tr$ of draw records $(c:s)$,
\[
  \seqt{\Gamma}{\Delta}{A}{\Tr} \;:=\;
  \seq{\Gamma}{\Delta,\trz{\Tr}}{A},
  \qquad
  \wt{\Tr} := \textstyle\prod_{(c:s)\in\Tr}\rho(c).
\]
\end{definition}

Every piece of trace bookkeeping reduces to context bookkeeping already
present in the base calculus:

\begin{center}
\begin{tabular}{ll}
\toprule
trace-zone behavior & context behavior it reduces to\\
\midrule
cut concatenates traces & cut merges linear contexts\\
additive lockstep ($\addconj$R shares $\Tr$) & $\addconj$R shares $\Delta$ (verbatim)\\
promotion is trace-empty & $\bang$R requires an empty linear zone (verbatim)\\
ghost pads $\Tr$ & weakening, restricted to tokens\\
weight multiplies along cut & $\wt{\Tr_1\uplus\Tr_2}=\wt{\Tr_1}\cdot \wt{\Tr_2}$\\
\bottomrule
\end{tabular}
\end{center}

\subsection{The rules}
\label{subsec:rules}

Three rules are added to the base calculus:

\begin{mathpar}
\inferrule*[right=$\exists_\rho$R$(c)$]
  {\seq{\Gamma}{\Delta}{A[c/x]}}
  {\seq{\Gamma}{\Delta,\dtok{c}{s}}{\erho{x}{s}A}}
\and
\inferrule*[right=$\exists$L]
  {\seq{\Gamma}{\Delta, A[a/x]}{C} \\ (a\ \text{fresh})}
  {\seq{\Gamma}{\Delta, \erho{x}{s}A}{C}}
\and
\inferrule*[right=ghost]
  {\seq{\Gamma}{\Delta}{C}}
  {\seq{\Gamma}{\Delta,\dtok{c}{s}}{C}}
\end{mathpar}

$\exists_\rho$R \emph{consumes} a token: given that a draw of $c$ at
sort $s$ occurred, assert the weighted existential via $c$.
$\exists$L is the standard eigenvariable rule, token-neutral.
$\ghost$ is weakening restricted to the token class: $\Delta$ proper
stays linear, the token sub-zone is affine (mixed-discipline zones are
standard).  No rule contracts a token; no rule proves one.

\begin{remark}[Why weight is not a grade]\label{rem:notagrade}
In the annotated design the principal reduction of
$\exists_\rho$R$(c)$ against $\exists$L substitutes $c$ for the
eigenvariable and deletes the $\exists_\rho$R node --- the only place
$\rho(c)$ was multiplied in.  On a two-member sort with
$\rho \equiv 1$ per member this rewrites $\langle 2\rangle$ to
$\langle 1\rangle$: annotated weight is not a proof invariant.  In
\will{} the same reduction \emph{keeps the token} (case (a) below); the
factor never lived on the deleted node, it lives in the endsequent.
\end{remark}

%%====================================================================
\section{Cut admissibility}
\label{sec:cut}
%%====================================================================

\begin{theorem}[Cut admissibility with exact trace conservation]
\label{thm:cut}
The three cuts of the base calculus (linear, lax, persistent) remain
admissible in \gill{} extended by the rules of \S\ref{subsec:rules},
with the same scope as the base theorem (the cut-free calculus excludes
the two premise-free oracle rules, which certificate elaboration
removes before verification).  Every reduction step preserves the
endsequent; hence the trace, and its weight, are conserved
\emph{exactly}, with raw $\QQ_{\geq0}$ priors.
\end{theorem}

\begin{proof}
By the base calculus's lexicographic induction: cut-formula weight
(with the first-order convention --- witnesses are terms, sort members
or constructor heads, never formulas, so witness substitution preserves
formula weight), then cut kind, then premise heights.  Endsequent
preservation holds by construction in every case; since the trace is
part of the endsequent, conservation is definitional.  The burden is
that the required reductions \emph{exist}.  We list the new or
newly-audited cases; all commutative cases not mentioned are those of
the base proof with tokens following their contexts.

\medskip\noindent\textbf{(a) Principal $\exists_\rho$R / $\exists$L.}
The crash site of the annotated design, resolved:
\[
\begin{array}{ll}
\pi' : \seq{\Gamma}{\Delta_1}{A[c/x]}
& \sigma' : \seq{\Gamma}{\Delta_2, A[a/x]}{C} \quad (a\ \text{fresh})\\[2pt]
\pi : \seq{\Gamma}{\Delta_1, \dtok{c}{s}}{\erho{x}{s}A}
& \sigma : \seq{\Gamma}{\Delta_2, \erho{x}{s}A}{C}\\[4pt]
\multicolumn{2}{l}{\mathrm{cut}(\pi,\sigma) :
  \seq{\Gamma}{\Delta_1, \dtok{c}{s}, \Delta_2}{C}}
\end{array}
\]
Reduct: substitute ($\sigma'[c/a]$ --- legal: $a$ fresh, $c$ ground,
tokens in $\Delta_2$ untouched by the fences), cut on the smaller
formula, and restore the token by one $\ghost$:
\[
  \mathrm{cut}(\pi', \sigma'[c/a]) :
  \seq{\Gamma}{\Delta_1,\Delta_2}{C}
  \;\;\leadsto\;\;
  \ghost :
  \seq{\Gamma}{\Delta_1,\Delta_2,\dtok{c}{s}}{C},
\]
the same endsequent.  The reduction \emph{needs} $\ghost$: it converts
a used draw into an unused one.  A system whose tokens must be exactly
consumed is not closed under cut reduction --- that is the annotated
design's crash restated as a general law, and why $\ghost$ is a rule,
not a leak.

\medskip\noindent\textbf{(b) Ghost commutation.}  $\ghost$ as the last
rule of either cut premise commutes below the cut (standard weakening
commutation; premise height strictly decreases).  $\ghost$ never
introduces the cut formula --- tokens are never cut formulas by (d) ---
so no principal case arises.

\medskip\noindent\textbf{(c) Identity cuts.}  id exists at all
formulas, in particular at $\erho{x}{s}A$ and at $\dtok{c}{s}$; a cut
against id reduces by deletion on either side, as always.

\medskip\noindent\textbf{(d) Cuts on tokens.}  The only derivations of
$\seq{\Gamma}{\Delta}{\dtok{c}{s}}$ end in id --- no right rule, no
checker conclusion (\S\ref{sec:calculus} fences).  Case (c) covers
them; no other principal case on tokens exists.

\medskip\noindent\textbf{(e) Duplication audit.}  The one duplication
mechanism is the persistent cut: $\mathrm{cut}^{\bang}$ substitutes a
promoted derivation at each copy.  Promotion requires the linear zone
empty --- which now includes tokens.  The duplicated derivation
therefore contains no draws: \textbf{no draw is ever duplicated}.  (The
counted exponential $\gbang{k}$ peels by \emph{splitting} context,
never by duplicating; harmless.)  The consequence deserves its name:
\textbf{no promotion through a draw} --- a derivation that draws is not
promotable.  Operationally: draws are linear, event-world phenomena;
the persistent world is deterministically re-derivable knowledge.  One
may bang the consequences of an outcome, never the luck itself.

\medskip\noindent\textbf{(f) Erasure audit.}  Reductions that delete a
subderivation: (i) principal $\addconj$ ---
$\mathrm{cut}(\addconj\text{R}(\pi_1,\pi_2), \addconj\text{L}_i(\sigma))$
keeps $\pi_i$ and erases the other; $\addconj$R's premises share one
$\Delta$, tokens included, so the endsequent loses nothing.
(ii) persistent weakening --- an unused persistent hypothesis erases
its promotion $\pi$; $\pi$'s linear zone is empty by (e), so no token
is lost.  (When $\adddisj$/$0$ left rules are added to the fragment,
$\adddisj$L shares $\Delta$ like $\addconj$R and $0$L erases into an
arbitrary endsequent, vacuously safe.)  \textbf{No draw is ever
erased.}

\medskip\noindent\textbf{(f$'$) The counted-promotion case.}  Found
fatal-as-stated by an adversarial audit of this proof.  As originally
written, the count-zero rule concluded $\gbang{K}A$ (with
$K\!=\!0$) from an \emph{arbitrary} linear context $D$ --- an affine
leak in the base calculus (it derives $a \vdash \gbang{0}b$,
contradicting $\gbang0 A \dashv\vdash 1$), harmless there because a
stranded linear atom leaves the derivation incompletable elsewhere,
but \emph{reachable} exactly when affine tokens exist: with
$D = \trz{\Tr}$, the principal cut against the count-zero left rule
reduces to that rule's premise and erases the tokens --- trace
conservation violated.  Two sound repairs exist: ghost-pad the reduct
by $|\Tr|$ (legal, since only tokens can inhabit that $D$ in a
completable derivation), or fix the rule to require an empty linear
zone like $1$R.  We took the root fix; with it the case is vacuous and
(f) stands.  The audit trail matters: the flaw was in the \emph{base}
calculus and invisible there; the affine token discipline is what made
it observable.

\medskip\noindent\textbf{(g) Lax and transport cuts.}  The graded
monad and transport rules are single-premise rules with persistent
theory premises (arithmetic side conditions).  Theory premises are
persistent-zone derivations --- clause resolution, no linear zone,
hence no tokens; witness substitution $[c/a]$ reaches them only as
ground-for-ground replacement in already-ground arithmetic goals
(tokens and grade terms are ground), preserving derivability.  Grades
compose exactly as in the base cut proof; tokens ride the linear
contexts orthogonally.  Delay grades and weight tokens never share a
slot: the ``product of two semirings on one judgment'' question
dissolves --- \textbf{grades on formulas, draws in the zone}.

\medskip\noindent\textbf{(h) Additive lockstep, derived.}
$\addconj$R gives both premises the same $\Delta$, hence the same
tokens; a branch that draws less discharges its surplus by $\ghost$.
The lockstep discipline that a standalone trace zone would have to
stipulate is here a consequence of context sharing.

\medskip\noindent\textbf{(i) Remaining commutative cases.}  Cut
formula non-principal on one side: push the cut past the last rule,
standard for every rule of the fragment.  The checks worth naming:
context-splitting rules ($\tensor$R, $\loli$L, the counted-peel and
transport rules) --- the cut lands in the branch holding the cut
formula, tokens follow their contexts; the stamp-retiming axiom ---
its theory premise is ground under witness substitution as in (g);
copy --- duplicates from $\Gamma$ only, token-free by (e).  The
induction measure is the base paper's, with $\ghost$ commutation
strictly decreasing premise height.
\end{proof}

%%====================================================================
\section{No-cloning: identity expansion fails at \texorpdfstring{$\exists_\rho$}{the weighted existential}}
\label{sec:nocloning}
%%====================================================================

\begin{theorem}[Identity-expansion failure]\label{thm:nocloning}
Identity expansion --- the derivation of $A \vdash A$ at compound $A$
from identities at atoms --- fails at $\exists_\rho$: there is no
token-free eta-expansion of
$\erho{x}{s}B \vdash \erho{x}{s}B$.
\end{theorem}

\begin{proof}
Expanding by $\exists$L opens a fresh eigenvariable $a$; concluding the
right side needs $\exists_\rho$R$(c)$ for some member $c$ --- which
consumes a token the sequent does not have.  By the fences no rule
proves a token, and $\ghost$ only discards.  With
Theorem~\ref{thm:cut}(d), every derivation of the sequent ends in id.
\end{proof}

Identity stays \emph{primitive} at $\exists_\rho$ (as it already is in
the base calculus, where id is general).  The physical reading: a
committed draw can be passed along whole, but not deconstructed and
re-derived --- re-derivation would be a second draw, visible as a
different endsequent with one more token.  A no-cloning-flavored
theorem at exactly the connective that carries randomness, in a
classical calculus, for free.  Its design consequence surfaces in the
focused system (\S\ref{sec:focused}): where standard focalization
arguments eta-expand identities at compound formulas, $\exists_\rho$
has no expansion, so the focused system must admit id at
$\exists_\rho$ as at an atom --- $\exists_\rho$ is a \emph{synthetic
atom on the identity dimension}.

%%====================================================================
\section{The ghost discipline}
\label{sec:ghost}
%%====================================================================

Call a derivation \emph{ghost-free} if every token in its endsequent is
consumed by an $\exists_\rho$R instance.

\begin{lemma}[Ghost commutation]\label{lem:ghostcomm}
$\ghost$ permutes below every rule of the calculus.  For
single-premise rules the permutation is verbatim.  For
context-splitting rules ($\tensor$R, $\loli$L, counted peel) the
parent's split is adjusted: the branch that held the ghosted token
proceeds with the smaller context, which its subderivation proves by
the induction hypothesis.  For $\addconj$R the token is ghosted in both
branches (shared context).  For $\exists_\rho$R the interaction is a
\emph{merge}, not a permutation: a $\ghost$ introducing $\dtok{c}{s}$
immediately consumed by $\exists_\rho$R$(c)$ cancels against it --- the
pair rewrites to the premise re-tokened from the ambient context.
\end{lemma}

\begin{proof}
Standard weakening commutation, by cases on the rule below the
$\ghost$; the merge case is immediate from the rule shapes.  Each move
either preserves the number of $\ghost$ occurrences or (merge)
decreases it; premise heights strictly decrease under each
permutation.
\end{proof}

\begin{lemma}[Stripping and padding]\label{lem:strip}
(i) Every derivation of $\seq{\Gamma}{\Delta,\trz\Tr}{A}$ contains a
ghost-free derivation of $\seq{\Gamma}{\Delta,\trz{\Tr_0}}{A}$ for some
$\Tr_0 \subseteq \Tr$: permute each $\ghost$ to the root
(Lemma~\ref{lem:ghostcomm}) and delete it there; induction on the
$\ghost$ count.
(ii) Conversely $\seq{\Gamma}{\Delta,\trz{\Tr_0}}{A}$ implies
$\seq{\Gamma}{\Delta,\trz{\Tr}}{A}$ for every $\Tr \supseteq \Tr_0$,
by applying $\ghost$.
\end{lemma}

\begin{corollary}[The counting class]\label{cor:countingclass}
Provable traces form the upward closure of ghost-free traces;
$\subseteq$-minimal provable traces are ghost-free.  The converse
fails: ghost-free is \emph{strictly} larger than minimal.  For
internal choice on the left,
\[
  (p\,c) \addconj (\erho{x}{s}\,p\,x) \;\vdash\; \erho{x}{s}\,p\,x
\]
is provable ghost-free at trace $\{c\}$ (project left, then
$\exists_\rho$R$(c)$) \emph{and} at the empty trace (project right,
then id --- the synthetic-atom identity of
Theorem~\ref{thm:nocloning}); the same separation arises from
$\adddisj$ on the right.  Ghost-free, not minimal, is the correct
counting class: both proofs carry mass.
\end{corollary}

\begin{remark}[Divergence audit]
Counting \emph{all} derivations would attach the factor
$\sum_k \big(\sum_c \rho(c)\big)^k = \infty$ to each proof through its
padded variants --- the trace-faithfulness objection of trace
types~\cite{LewPOPL2020} in proof-theoretic form.  Counting ghost-free
derivations excludes every padded variant.  Note that ghost marks live
on \emph{rule occurrences}, not in the trace: marked trace entries
would break Theorem~\ref{thm:cut}'s endsequent-on-the-nose
conservation, since reduction (a) creates a $\ghost$ occurrence while
preserving the endsequent.  This is exactly weakening's profile:
present in derivations, absent from canonical forms.
\end{remark}

\begin{proposition}[Conservativity]\label{prop:conserv}
Erasing tokens and mapping $\exists_\rho$R/$\exists$L to the standard
$\exists$R/$\exists$L sends \will{}-derivations to ILL derivations.
Conversely any ILL derivation lifts: take $\Tr$ = the multiset of
$\exists$R witnesses at the $\exists_\rho$ positions; each
subderivation's tokens are the witnesses of its own $\exists_\rho$R
occurrences, so context-splitting rules split $\Tr$ canonically
(disjoint subtrees) and context-sharing rules are balanced by
$\ghost$-padding each branch with the other's tokens.  Hence
$\seqt{\Gamma}{\Delta}{A}{\Tr}$ holds for some $\Tr$ iff the plain
judgment is derivable: the trace judgment is a conservative extension.
\end{proposition}

%%====================================================================
\section{The focused system}
\label{sec:focused}
%%====================================================================

The counting fragment is the multiplicative-additive core plus
$\exists_\rho$ and tokens:
\[
  P ::= a \mid \dtok{c}{s} \mid 1 \mid A \tensor B \mid \erho{x}{s}A
  \qquad
  N ::= A \loli B \mid A \addconj B
\]
with atoms positive.  ($\adddisj$ extends the fragment routinely; the
graded template rules --- monad, exponentials, transport, stamps ---
and the oracle rules stay outside: the boundary judgments of
\S\ref{sec:bridge} live in the collapse fragment, matching the
driver.)  The persistent zone is passive throughout (copy + id;
token-free by Theorem~\ref{thm:cut}(e)).

\subsection{Polarities and phases}

We use a standard two-phase focused presentation for intuitionistic
linear logic~\cite{Andreoli1992,ChaudhuriPfenningPrice2008}:
\emph{inversion} eagerly decomposes negatives on the right and
positives on the left; at a \emph{stable} sequent a focus is chosen
(right focus on a positive goal, left focus on a negative hypothesis);
focus persists through the chosen formula's spine.  The additions:

\begin{itemize}
\item \textbf{$\dtok{c}{s}$ is a left-passive positive atom.}  It has
  no left decomposition; during left inversion it moves to the stable
  context like an atom.  It is consumed only as the side condition of
  the focused $\exists_\rho$R:
  \[
  \inferrule*[right={$\exists_\rho$R$(c)$, focused}]
    {\Gamma;\;\Delta \focR A[c/x]}
    {\Gamma;\;\Delta, \dtok{c}{s} \focR \erho{x}{s}A}
  \]
  The token is consumed \emph{in phase} and focus continues on
  $A[c/x]$.  Choosing $c$ is choosing which token to spend --- exactly
  the non-invertible choice focusing isolates into the focus phase.
\item \textbf{$\exists$L in inversion, or suspension.}  Left inversion
  on $\erho{x}{s}A$ either applies the (token-neutral, eigenvariable)
  $\exists$L, or \emph{suspends} the formula to the stable context as a
  synthetic atom, to be consumed later only by id against a
  right-focus occurrence of the same formula.  Suspension is forced by
  Theorem~\ref{thm:nocloning}: id at $\exists_\rho$ has no expansion,
  so a derivation that passes a wave along whole cannot be represented
  by decompose-and-rebuild.  (For the boundary judgments of
  \S\ref{sec:bridge} the goal contains no $\exists_\rho$, id at a
  suspended wave can never fire, and suspension is vacuous ---
  decomposition is then exhaustive and the phase discipline fully
  canonical.)
\item \textbf{No ghost.}  The focused system has no $\ghost$ rule:
  ghost-freeness is not a side condition on the count but a structural
  property of the system --- weakening's profile under focusing,
  exactly (\S\ref{sec:ghost}).
\item \textbf{Identity.}  id closes a right focus on an atom, a token,
  or a (suspended) $\exists_\rho$ formula against an identical stable
  hypothesis.
\end{itemize}

\subsection{Completeness}

\begin{theorem}[Focalization]\label{thm:focal}
Every ghost-free cut-free derivation of a counting-fragment sequent
has a focused derivation of the same endsequent, and conversely
(soundness) every focused derivation erases to a ghost-free cut-free
derivation.
\end{theorem}

Soundness is immediate: erase phase markers; each focused rule is an
instance of an unfocused rule; no rule erases to $\ghost$.  For
completeness we follow the standard permutation
strategy~\cite{Andreoli1992,Laurent2004,MillerSaurin2007,Simmons2014}:
(1) invertible rules permute with each other and can be applied
eagerly; (2) each non-invertible rule permutes downward until it either
meets its phase or reaches a stable sequent, and chains of positive
rules on one formula's spine can be made contiguous (focus
persistence).  The base MALL cases are textbook; we work every case our
additions create.  Throughout, ``preserves ghost-freeness'' is
immediate for each move because no permutation square below introduces
a $\ghost$ occurrence.

\begin{lemma}[Token passivity]\label{lem:passivity}
In the counting fragment, the only rules that mention tokens are
$\exists_\rho$R (consumes one) and id (at a token).  In particular no
invertible rule consumes, produces, or decomposes a token, and a token
in $\Delta$ commutes inertly with every inversion step.
\end{lemma}

\begin{proof}
By inspection of the rule set; $\ghost$ is absent from ghost-free
derivations by definition, and the fences exclude tokens from every
program-defined rule.
\end{proof}

\begin{lemma}[Principal occurrences of a left wave]\label{lem:leftwave}
In a cut-free derivation, each occurrence of $\erho{x}{s}A$ in a linear
context is principal at exactly one rule: $\exists$L or id.
\end{lemma}

\begin{proof}
Linearity: the occurrence is consumed exactly once; the only rules
that can consume a left $\exists_\rho$ are $\exists$L and id (no other
left rule matches it, and $\ghost$ only consumes tokens).
\end{proof}

Lemma~\ref{lem:leftwave} is what justifies the suspension discipline:
trace the occurrence upward; if it is principal at $\exists$L, permute
that instance down into the inversion phase (the standard invertible
permutations, token-inert by Lemma~\ref{lem:passivity}); if principal
at id, suspend it --- it behaves as an atom along its entire path.

\begin{lemma}[$\exists_\rho$R permutations]\label{lem:erhoperm}
Let a ghost-free cut-free derivation contain an $\exists_\rho$R$(c)$
instance immediately above a rule $R$ that is not part of its focus
phase.  Then the two permute (the $\exists_\rho$R moves below $R$, or
$R$ moves above it) with the same endsequent, no new $\ghost$, and all
token threading determined by the contexts.  The cases:

\emph{(i) $R$ invertible on the right} ($\loli$R, $\addconj$R): the
goal of $\exists_\rho$R is the whole succedent, so $R$'s principal
formula is the same $\erho{x}{s}A$ --- impossible ($\exists_\rho$ is
not negative); no interaction.

\emph{(ii) $R$ invertible on the left} ($\tensor$L, $1$L,
$\exists$L on a different formula): $R$ does not touch the token
(Lemma~\ref{lem:passivity}) nor the succedent; permute $R$ below:
\[
\inferrule*{\inferrule*{\seq{\Gamma}{\Delta', A[c/x]}{G}}
            {\seq{\Gamma}{\Delta, A[c/x]}{G}}}
           {\seq{\Gamma}{\Delta, \dtok{c}{s}}{\erho{x}{s}G'}}
\;\;\leadsto\;\;
\inferrule*{\inferrule*{\seq{\Gamma}{\Delta', A[c/x]}{G}}
            {\seq{\Gamma}{\Delta', \dtok{c}{s}}{\erho{x}{s}G'}}}
           {\seq{\Gamma}{\Delta, \dtok{c}{s}}{\erho{x}{s}G'}}
\]
(eigenvariable side conditions unaffected: $c$ is ground).

\emph{(iii) $R = \tensor$R or $\loli$L (context split), with the
$\exists_\rho$R in one branch}: the token lives in that branch's
context; moving the boundary between phase and stable sequent
re-associates the split so that the token stays with its branch ---
the standard multiplicative permutation with one extra formula tracked,
possible because the token is not decomposed by any rule the split
could force (Lemma~\ref{lem:passivity}).

\emph{(iv) $R = \addconj$R below}: both branches share
$\Delta, \dtok{c}{s}$.  In a ghost-free derivation \emph{each} branch
consumes the token by its own $\exists_\rho$R (there is no other
consumer), so the permutation squares apply branchwise; no
$\ghost$-balancing is ever needed --- this is
Theorem~\ref{thm:cut}(h) restricted to the ghost-free class.

\emph{(v) $R = $ id}: id has no premises; vacuous.
\end{lemma}

\begin{proof}
Each square is checked by writing out both composites; the contexts
match multiset-theoretically because the token is a formula like any
other, and no square introduces $\ghost$.
\end{proof}

\begin{proof}[Proof of Theorem~\ref{thm:focal} (completeness)]
Induction on the derivation with the standard measure (number of rule
instances not in phase position).  Apply invertible permutations to
make the inversion phase maximal at the root --- possible by the base
invertibility lemmas extended to $\exists$L via
Lemma~\ref{lem:leftwave}: each left wave is decomposed eagerly, or
suspended if its principal rule is id (then it is atom-like for the
rest of the argument).  At the resulting stable sequent, pick the rule
that begins the derivation's next non-invertible cluster; by
Lemmas~\ref{lem:passivity} and~\ref{lem:erhoperm} (and the base
permutations for the MALL connectives) its entire positive spine can
be made contiguous --- in particular an $\exists_\rho$R$(c)$ instance
permutes down to join the right-focus phase of its goal formula, its
token consumed in phase.  Where the standard argument eta-expands an
identity at a compound formula, id at $\exists_\rho$ (and at tokens) is
instead kept as a focused leaf --- legitimate because $\exists_\rho$ is
synthetic (Theorem~\ref{thm:nocloning} means no expansion exists to
normalize it away, and none is needed).  Each step reduces the
measure; the process terminates in a focused derivation of the same
endsequent, ghost-free at every intermediate stage since no move
introduced a $\ghost$.
\end{proof}

\begin{definition}[The count]\label{def:count}
For a counting-fragment sequent whose goal contains no
$\exists_\rho$,
\[
  N(A,\Tr) \;:=\;
  \#\{\text{focused, cut-free, ghost-free derivations of }
  \seq{\Gamma}{\Delta,\trz\Tr}{A}\}.
\]
By the scope condition suspension is vacuous
(\S\ref{sec:focused}), so the focused system is fully canonical and
$N$ is well-defined.  Theorem~\ref{thm:focal} makes $N$ count exactly
the ghost-free cut-free proofs up to phase-permutation.
\end{definition}

%%====================================================================
\section{Universality: the proof-counting provenance polynomial}
\label{sec:bridge}
%%====================================================================

\begin{definition}\label{def:poly}
For a boundary judgment ($\Delta$, $A$ token-free), the
\emph{proof-counting provenance polynomial} and the $K$-mass are
\[
  P_A(X) := \sum_{\Tr} N(A,\Tr)\cdot X^{\Tr} \;\in\; \NN[X],
  \qquad
  \mu_K(A) := \mathrm{eval}_K(P_A)\ \text{at}\ X_c \mapsto \rho(c),
\]
for any commutative semiring $K$; a trace \emph{is} a monomial.
\end{definition}

By $\NN[X]$-initiality~\cite{GreenKT2007} every semiring-valued reading
factors through $P_A$.  Because tokens are linear, each derivation
contributes exactly one monomial --- its endsequent's token multiset:
where the provenance-semiring literature sums monomials over proof
trees, \will{} reads the monomial off the sequent.  The instances:

\begin{itemize}
\item $K = (\QQ_{\geq0},+,\cdot)$ is \will{}'s probabilistic mass.
  The identification of $\mu_{\QQ_{\geq0}}$ with the measure computed
  by the engine's exact mode, and the convergence (least-fixpoint
  masses for recursive sorts under Chi--Geman
  subcriticality~\cite{ChiGeman1998,EtessamiYannakakis2009}) and
  sampler-unbiasedness statements, are the companion theory's
  T1--T3~\cite{THY0026}; \S\ref{sec:adequacy} supplies the missing
  proof-theoretic leg.
\item $K = (\min,+)$ is the cost optimum of the graded sibling
  calculus: idempotence collapses the sum onto the single best
  derivation, which is exactly why \gill{} needs no trace and a
  per-proof grade suffices.  The gap between the two calculi is the
  ambiguity hierarchy of weighted automata in proof-theoretic form.
\item $K = \CC$ is reserved.  Evaluation is still a homomorphism, but
  interference (cancellation) is forgetting that no trace-level
  statement can undo: the trace level is the initial object, each
  quantitative face a quotient of it.
\end{itemize}

Cut admissibility is proved once (Theorem~\ref{thm:cut}); each face
inherits invariance through its homomorphism.

%%====================================================================
\section{Driver adequacy}
\label{sec:adequacy}
%%====================================================================

\subsection{The decimation driver}

The operational semantics extends the timed committed-choice engine of
the companion paper~\cite{TillPaper}.  A rule consequent may contain
$\erho{x}{s}A$; firing such a rule \emph{suspends} the existential as a
linear \emph{wave} fact.  Plain execution never touches a wave (the
opt-in principle: unresolved choice is inert).  The \emph{decimation
driver} runs the loop
\[
  \text{settle} \to \text{open} \to \text{settle} \to \text{draw} \to
  \text{substitute} \to \cdots
\]
\emph{Open} replaces each suspended wave by its body instantiated with
a fresh existential variable (an $\exists$L instance, operationally);
\emph{settle} runs ordinary committed-choice firing to quiescence, so
that program rules may derive persistent \emph{bias} facts mentioning
the wave variable --- monotone conditioning; the \emph{posterior} of a
wave is $\rho(\cdot)$ times the product of its distinct bias factors;
\emph{draw} picks a wave (minimum posterior entropy first, a
semantics-free policy) and a member (pseudorandomly by exact rational
weights, or exhaustively in exact mode), substitutes the witness ---
for a constructor member, the head only, with argument positions
becoming new waves (lazy PCFG) --- and repeats until no waves remain.
A run's \emph{trace} is its multiset of drawn heads; its mass the
product of drawn posterior weights.

\subsection{Adequacy}

Fix a program whose settle segments are \emph{conflict-free}: between
draws, committed-choice firing is deterministic up to the enumerated
weighted-choice forks (the engine enforces this loudly in exact mode
--- a genuine scheduling conflict is adversarial nondeterminism,
outside the measure).  Fix a draw policy (a deterministic choice of
which wave to draw next; the implementation's entropy policy is one).
A \emph{collapse path} is a maximal driver run; its \emph{trace} the
multiset of drawn heads with their sorts.

\begin{theorem}[Adequacy]\label{thm:adequacy}
Let $S_0$ be an initial boundary state and consider the boundary
judgments $\Gamma;\enc{S_0},\trz{\Tr} \vdash \laxb{\enc{F}}$ for final
ground states $F$ (where $\enc{\cdot}$ is the standard state encoding
and the persistent zone $\Gamma$ holds the program's derived
knowledge).  The map that elaborates a collapse path to a derivation
--- each settle segment to its (cited) exact fire chain, each wave's
open+draw pair to the derived inference
$\mathrm{cut}(\exists_\rho\text{R}(c), \exists\text{L})$ at the
opening position --- is a bijection between fixed-policy collapse
paths from $S_0$ with trace $\Tr$ and final state $F$, and the focused
cut-free ghost-free derivations counted by $N(\laxb{\enc F},\Tr)$
over the collapse fragment.
\end{theorem}

\begin{proof}
\emph{Well-defined.}  The elaboration is total on driver runs: settle
segments elaborate exactly (the base engine's trace-elaboration
theorem~\cite{TillPaper}); the open+draw pair is the principal-cut
reduct of Theorem~\ref{thm:cut}(a) read backwards --- the driver's
substitution of the drawn witness is the reduct's substitution, and
the token the certificate keeps is the endsequent record of the
consumed $\exists_\rho$R.  Post-hoc grounding (substituting each final
witness through the whole recorded run) is sound because bias
derivation is monotone --- facts derivable between an opening and its
draw remain derivable after --- and existential variables are opaque
constants to matching, so a grounded instance of a legal firing is a
legal firing.  Ghost-freeness: every minted token is consumed by its
own draw; unopened or unobserved waves mint no token and contribute no
factor.

\emph{Injective.}  Two distinct fixed-policy paths first differ at a
draw (segments are deterministic by conflict-freedom; the policy
determines which wave is drawn next).  A differing member gives a
different token/witness at the corresponding $\exists_\rho$R; a
differing outcome of the same member changes nothing (draws are
member-level).  Hence the elaborated derivations differ.

\emph{Surjective.}  Given a focused ghost-free derivation $\pi$ of the
boundary judgment, read off its $\exists_\rho$R instances.  Because
the goal contains no $\exists_\rho$, every left wave in $\pi$ is
decomposed (suspension is vacuous, Definition~\ref{def:count}), and
the derivation tree's structure is the org chart, not a schedule: the
draws are partially ordered by subterm dependency only.  The policy
linearizes this partial order in exactly one way compatible with the
driver's segments (which are determined by quiescence), yielding a
path whose elaboration is $\pi$: the fire chains match by the base
adequacy theorem, the draws by construction.

The two directions compose to identities because elaboration forgets
only the linearization the policy reconstructs.
\end{proof}

\begin{corollary}[Policy independence and the mass identity]
\label{cor:mass}
$N(\cdot)$, and hence the mass
$\mu_{\QQ_{\geq0}}(\laxb{\enc F}) = \sum_\Tr N\cdot\rho^\Tr$, is
independent of the draw policy: each derivation has exactly one
linearization per policy.  Exact mode's enumeration computes this mass
(it enumerates fixed-policy paths); the importance-weighted sampler
targets it (unbiasedness per attempt is~\cite{THY0026}'s T3).
\end{corollary}

\begin{remark}
The scope conditions earn their keep operationally.  Conflict-freedom
is checked, not assumed: the exact-mode engine enumerates weighted
forks and \emph{fails loudly} on a genuine scheduling conflict.  The
boundary-judgment restriction (goal without $\exists_\rho$) is the
statement that runs end ground; suspended waves surviving to the
boundary are precisely the case the certificate handles by
$\exists$-closing the goal (\S\ref{sec:impl}).
\end{remark}

%%====================================================================
\section{Implementation: certificates that carry their trace}
\label{sec:impl}
%%====================================================================

\will{} is implemented in the \calc{} sandbox on top of the engine of
the companion paper.  We describe what is new, and what the theorems
of this paper buy operationally.

\subsection{The rules as data, the checker as judge}

The calculus definition is declarative (a \texttt{.rules} file); the
three rules of \S\ref{subsec:rules} are template rules with binder
opening --- $\exists_\rho$R is mechanically a \emph{left} rule on the
token whose succedent pattern forces the token's sort to match the
wave's, with the witness bound from the token; $\ghost$ is marked
affine and realized in proof search as a last-resort focus choice plus
boundary discharge (root leftovers and additive-branch balancing ---
Theorem~\ref{thm:cut}(h) operationalized); a greedy searched ghost
steals tokens under committed premise search and a lazy one starves
additive branches, so the boundaries are the only sound insertion
points.  Search-found derivations are re-verified by an LF-style
kernel that re-derives every step from the rule data.

One step is checker-bound rather than searched: \texttt{@draw}, the
derived rule $\mathrm{cut}(\exists_\rho\text{R},\exists\text{L})$ that
records one collapse event.  Its checker re-derives everything from
program data --- classifier membership, witness-head discipline, the
declared prior against the recorded weight, and the opened-body
context arithmetic --- and is the \emph{only} place tokens are minted:
programs can neither produce nor match \textsf{drawn} (a load-time
fence), so the fences of \S\ref{sec:calculus} are enforced, not
conventions.

\subsection{Certified runs}

The central artifact is \texttt{certifyCollapse}: a complete
stochastic run elaborates into \emph{one} kernel-checked derivation
\[
  \Gamma;\;\Delta_0, \trz{\Tr} \;\vdash\; \laxb{\textstyle\bigotimes
  \text{residual}}
\]
with one \texttt{@fire} node per firing (fully re-derived; the base
paper's certified-execution pipeline), one \texttt{@draw} node per
wave at its \emph{opening} position with the post-hoc-grounded witness
--- a rung-2 witness tree certifies as the composite of its per-head
draws, one token consumed per drawn head; an existential variable in
the witness is a wave dropped unobserved: no choice, no token, no
factor --- and token-free \emph{open} nodes ($\exists$L with
syntactically-checked eigenvariable freshness) for suspensions that
were never drawn.  Skolems surviving to the residual are
$\exists$-closed in the goal.  The endsequent carries $\trz\Tr$:
\textbf{the run's prior mass is $\prod_{\Tr}\rho(c)$, computable from
the certificate's endsequent alone} --- Definition~\ref{def:trace}
made executable.  Tampering is caught: a doctored draw weight fails
the checker against the declared prior; removing a token from the
endsequent makes the draw nodes unverifiable.

\subsection{Case study: wave function collapse}

The demo program is a four-cell, three-tile (\textsf{sea},
\textsf{coast}, \textsf{land}; priors $2{:}1{:}2$) strip with the
constraint that sea and land are never adjacent.  Each cell's wave is
one $\exists_\rho$; the constraint is two \textsf{forbid} facts and
\emph{one} bias rule (a collapsed cell biases its neighbour's wave
variable to zero) --- no domain bitmasks, no hand-compiled per-mask
weighted choices, which is what the same demo required before the
weighted existential existed.  The propagate-then-collapse order of
classical WFC is not programmed: it \emph{is} the minimum-entropy draw
policy, since a biased wave has lower posterior entropy than an open
one.  Exact mode enumerates the constrained measure: 41 of $3^4=81$
beaches carry mass, total $217$ of the unconstrained $5^4 = 625$ ---
matching brute-force enumeration exactly.  Every sampled run
certifies; the certificates' token products equal the reported masses
(bias factors are conditioning, derived inside the certificates' fire
segments; the endsequent weight is the prior product, per
Definition~\ref{def:trace}).

\subsection{Executable metatheory}

The theorems of this paper are pinned by the test suite: the
provability and refutation patterns of
\S\S\ref{sec:cut}--\ref{sec:ghost} (token discipline, ghost boundaries,
additive balancing, no-promotion-through-a-draw, the count-zero audit
repair, eigenvariable soundness) run as kernel-verified derivations;
the driver, the checker, and the certificates are exercised across
unbiased, biased, recursive (PCFG), skolem, and correlated-wave
programs, with adversarial no-fast-path arms re-running the suites
with all arithmetic acceleration disabled.  The full fast suite is
$\sim$3.5k tests; the will-specific additions of this work are
$\sim$90.

%%====================================================================
\section{Related work}
\label{sec:related}
%%====================================================================

\paragraph{Weighted quantifiers and probabilistic proof theory.}
Riesz modal logic~\cite{LucasMio2022} is, to our knowledge, the only
prior \emph{syntactic} cut-elimination result for a quantitatively
valued judgment; it works over hypersequents with scalar-multiplied
contexts, and its combination with linear multiplicatives is open.
Counting propositional logic~\cite{AntonelliDLP2022} interprets
counting quantifiers by measure but establishes cut admissibility
semantically, not syntactically.  The algebraic
$\lambda$-calculus~\cite{Vaux2009} shows where scalar sums survive
reduction (positive coefficients) and where they collapse.  \will{}
sidesteps the sum-valued judgment entirely: sequents count, semirings
evaluate.  Bayesian proof-nets~\cite{DiGuardiaEF2024} put probability
tables inside proof-net boxes with structure-preserving cut
elimination --- probabilities in \emph{boxes}, not hypotheses in the
\emph{context}, and no weighted quantifier.

\paragraph{Randomness as a resource.}  Dahlqvist and
Kozen~\cite{DahlqvistKozen2020} treat random sources as linear
resources --- denotationally, in a semantics of higher-order
probabilistic programs, with no proof system, no token hypotheses, no
cut theorem.  Time credits~\cite{CharguerraudPottier2017} and error
credits~\cite{Eris2024} are the closest proof-\emph{system} mechanism:
consumable hypotheses whose presence bounds cost or error.  Credits
are anonymous and additive (spend $n$ of a homogeneous budget);
draw tokens are \emph{named} witnesses with \emph{multiplicative}
per-constructor priors, consumed by a quantifier rule --- the
endsequent does not bound the weight, it \emph{is} the weight.
Nominal logic's fresh names~\cite{Pitts2003,MillerTiu2005} are
non-duplicable for scope hygiene; they carry no weight and no
quantifier consumes them.

\paragraph{Probabilistic logic programming.}  Sato's distribution
semantics~\cite{Sato1995} and PRISM's \texttt{msw} switches are the
operational ancestors: a run is a deterministic derivation conditional
on a random record.  The record there is a persistent, memoized fact
table with no proof theory; here it is a linear zone with a cut
theorem, and the difference is what makes weight a proof invariant.
Trace types~\cite{LewPOPL2020} type the shape of a probabilistic
program's trace and are the source of the faithfulness objection that
forced the ghost discipline (\S\ref{sec:ghost}): a counting class that
admits padded traces diverges.

\paragraph{Provenance.}  Provenance
semirings~\cite{GreenKT2007,GradelTannen2024} sum $\NN[X]$-monomials
over derivation trees, with $\exists$ interpreted as semiring sum.
\will{} sharpens the connection at the object level: linearity of
tokens means each derivation \emph{is} one monomial, read off its
endsequent, and cut elimination --- absent from the database setting
--- preserves it.

\paragraph{Focusing.}  The focused system and its completeness
argument ride on the standard
line~\cite{Andreoli1992,Laurent2004,MillerSaurin2007,ChaudhuriPfenningPrice2008,Simmons2014}.
The additions --- a left-passive resource consumed as a focus side
condition, and a compound connective that must be treated as a
synthetic atom because identity expansion \emph{provably} fails at it
--- appear to be new as a package; synthetic atoms per se are
folklore~\cite{Zeilberger2008}.

\paragraph{Graded modalities.}  The graded-necessity line
(bounded linear logic~\cite{GirardScedrovScott1992}, graded
sequent calculi, Granule~\cite{OrchardLiepeltEades2019}) and the
graded-monad line~\cite{Katsumata2014} grade \emph{formulas}.  The
companion paper~\cite{TillPaper} put a delay grade on the lax modality
with cut admissibility; this paper's point is precisely that weight
should \emph{not} be another grade: grades on formulas, draws in the
zone (Theorem~\ref{thm:cut}(g)).

\paragraph{Quantum analogies.}  The identity-expansion failure
(\S\ref{sec:nocloning}) is a no-cloning statement at the proof level;
we use it as a design constraint, not a quantum claim.  The reserved
$\CC$ face of \S\ref{sec:bridge} records why a trace-level calculus
cannot see interference: evaluation homomorphisms only forget.

%%====================================================================
\section{Conclusion}
\label{sec:conclusion}
%%====================================================================

A measure-weighted existential fits linear logic once the weight stops
being an annotation and becomes a hypothesis.  The internalized trace
turns every hard question into an old one: trace conservation is
endsequent preservation; additive lockstep is context sharing; the
duplication and erasure dangers are the promotion condition; the
counting class is weakening's focusing story.  The three structural
theorems --- ghost forced by reduction, no promotion through a draw,
no-cloning at the weighted quantifier --- are not costs of the design
but its content, each matching an operational invariant of the
stochastic engine the calculus certifies.

Future work: mechanizing Theorem~\ref{thm:cut} (the internalization
makes the calculus small --- standard ILL plus three rules --- so the
graded base is the larger mechanization burden); the mass-judgment
stretch face (sequents proving total mass, with the Riesz-style price
tag); inside-mass conditioning and datasort-constrained domains for
the driver; and the $\CC$ face, which requires leaving the trace
quotient altogether.

\bibliographystyle{alpha}
\bibliography{refs}

\end{document}
